\begin{figure}[!t]
\centering
\begin{subfigure}[t]{0.323\linewidth}
\centering\includegraphics[width=\linewidth]{figures/pdf/figF_bert_baseline.pdf}
\caption{Subgroup contrasts at dose zero.}
\end{subfigure}\hfill
\begin{subfigure}[t]{0.323\linewidth}
\centering\includegraphics[width=\linewidth]{figures/pdf/figF_bert_mid.pdf}
\caption{Subgroup contrasts at dose $0.5$.}
\end{subfigure}\hfill
\begin{subfigure}[t]{0.323\linewidth}
\centering\includegraphics[width=\linewidth]{figures/pdf/figF_bert_controls.pdf}
\caption{Constructed target responses.}
\end{subfigure}
\caption{\textbf{Subgroup localization and exact-duplication checks.} Fingerprint coordinates are arithmetic differences of archived context means; lines in (c) connect the observed dose-grid values. No missing confidence interval is reconstructed. The clone residual is retained at its actual nonzero numerical value, although it is visually indistinguishable from zero on the common scale.}
\label{fig:legacy}
\end{figure}

\section{Additional sentiment-classifier audits}
\label{app:legacy}

The SST-2 experiments complement the main generative-model study by examining identifiable input subgroups and constructed target relations. They contain BERT, DistilBERT, RoBERTa, ALBERT, and XLNet classifiers under matched token masking. They are not described as a modern generative-model benchmark.

\paragraph{Context contrasts.}
For each model and dose, the horizontal fingerprint coordinate subtracts mean PIER on the reported medium-length bin (8--15) from that on the long-sentence bin (more than 15); the vertical coordinate subtracts the no-negation group from the explicit-negation group. Figure~\ref{fig:legacy}(a)--(b) uses the stored subgroup means at doses $0$ and $0.5$. ALBERT moves from positive contrasts at dose zero to negative contrasts at dose $0.5$. DistilBERT has a negative negation contrast at both doses but a positive length contrast. These are descriptive subgroup contrasts, not causal effects of sentence length or negation.

\paragraph{Constructed targets.}
A separate fixed-peer audit compares a DistilBERT clone, RoBERTa, and a fine-tuned DistilBERT. The clone's maximum residual is $1.31\times10^{-8}$ across the masking grid. RoBERTa and the fine-tuned target have nonzero gaps that generally rise with dose (Figure~\ref{fig:legacy}c). The fine-tuned target already has a residual of $0.120$ at baseline; stronger masking increases an already present gap. These controls establish different response relations without attributing the observed differences solely to architecture or a particular training mechanism.


